\chapter{Il formato dell'archivio}
\label{cap:formato}

\section{File e messaggi}
\label{sec:file}

L'archivio è fatto di un file GRIB2 per ogni istante di validità. Ogni file orario contiene 90
variabili in 246 messaggi GRIB, sempre gli stessi: l'elenco viene da un unico registro del
convertitore, e ogni file prodotto viene confrontato con quell'elenco, così un campo che non si è
potuto calcolare non passa inosservato. I nomi dei file seguono lo schema
\campo{ailam-an-cima-3km-{year}-{year}-1h-v1-{YYYYMMDD}{HH}.grib}, in cui l'anno compare due volte
con lo stesso valore e la parte finale è l'istante di validità.

I messaggi usano il template di griglia 3.10 (Mercator), il centro di origine \campo{ecmf} e la
compressione \campo{grid_ccsds} a 24 bit per valore. Ogni parametro porta il paramId e lo
shortName ECMWF \parencite{paramdb}. Gli shortName distinti sono 89 per 90 variabili, perché
\campo{z} compare due volte, come geopotenziale sui livelli di pressione e come geopotenziale di
superficie, cioè l'orografia. I campi definiti solo su una parte del dominio portano una bitmap
GRIB, e i punti fuori dal dominio di definizione sono punti mancanti, non zeri
(\cref{sec:mascheratura}).

Un file orario pesa in media \SI{737}{MB} sui 184 file di settembre 2024, e \SI{661}{MB} sugli 8
file delle 00 UTC fra questi; sui 476 file di febbraio e marzo 2024 la media è \SI{765}{MB}. Ne
risultano fra \SI{17.7}{GB} e \SI{18.4}{GB} al giorno e fra \SI{6.5}{TB} e \SI{6.7}{TB} all'anno.

\section{Geometria}
\label{sec:geometria}

La griglia dichiara la forma della Terra con \campo{shapeOfTheEarth=1} e raggio
\SI{6370000}{m}, la sfera del modello (\cref{sec:dominio}). Con questa dichiarazione le coordinate
che eccodes ricostruisce dalla definizione di griglia coincidono con quelle del modello entro
\SI{2.009}{m} su tutto il dominio. Senza la dichiarazione, la griglia ricostruita da eccodes si
sposta fino a circa \SI{1.1}{km} al bordo nord.

\section{Livelli verticali}
\label{sec:livelli}

Le 13 variabili tridimensionali stanno sugli stessi 13 livelli di pressione,
\SIlist{1000;925;850;700;600;500;400;300;250;200;150;100;50}{\hPa}. Sono i livelli richiesti dal
progetto LAM AI, un sottoinsieme dei 37 livelli di pressione di ERA5 \parencite{era5doc}. Degli
altri, gli 8 sopra \SI{50}{\hPa} stanno sopra il top del modello e non si possono produrre.

I 246 messaggi si ripartiscono fra i tipi di livello della \cref{tab:tipi-livello}.

\begin{table}[htbp]
  \centering
  \small
  \caption{Tipi di livello dei 246 messaggi di un file orario.}
  \label{tab:tipi-livello}
  \begin{tabularx}{\textwidth}{@{}lrX@{}}
    \toprule
    Tipo di livello & Messaggi & Contenuto \\
    \midrule
    \campo{isobaricInhPa}      & 169 & 13 variabili su 13 livelli di pressione \\
    \campo{heightAboveGround}  & 11  & vento a 10, 100 e \SI{200}{m}; temperatura, temperatura di rugiada e umidità relativa a \SI{2}{m}; \campo{10fg}; \campo{vwsh} \\
    \campo{surface}            & 47  & campi statici, suolo, neve, precipitazione, deflusso e radiazione in superficie \\
    \campo{entireAtmosphere}   & 10  & integrali di colonna e coperture nuvolose \\
    \campo{nominalTop}         & 5   & radiazione al top dell'atmosfera \\
    \campo{mostUnstableParcel} & 2   & \campo{mucape}, \campo{mucin} \\
    \campo{meanSea}            & 1   & \campo{msl} \\
    \campo{entireLake}         & 1   & \campo{dl} \\
    \bottomrule
  \end{tabularx}
\end{table}

\section{Convenzione temporale}
\label{sec:tempo}

I campi istantanei usano il template di prodotto 4.0, e il loro tempo di riferimento
(\campo{dataDate}, \campo{dataTime}) coincide con l'istante di validità. I 19 campi accumulati
sono la precipitazione totale (\campo{tp}), la pioggia (\campo{tirf}), la neve
(\campo{sf}), i tre deflussi (\campo{sro}, \campo{ssro}, \campo{ro}) e i 13 flussi radiativi in
superficie e al top dell'atmosfera. Usano il template 4.8, perché le definizioni ECMWF di questi
parametri richiedono la chiave \campo{typeOfStatisticalProcessing}, che nel template 4.0 non esiste
\parencite{eccodes}. Come in ERA5 \parencite{era5doc}, ogni campo accumulato copre l'ora
che termina all'istante di validità $H$; nell'archivio il tempo di riferimento è $H-1$ e
l'intervallo di accumulazione (\campo{stepRange}) è 0--1. Il valore è la differenza fra due ore successive della stessa run, per
esempio $\campo{tp}(H) = [\campo{RAINNC}(H) - \campo{RAINNC}(H-1)]/1000$, con \campo{RAINNC} in
\si{mm} e \campo{tp} in \si{m} \parencite{issue70}. Il ventesimo campo con template 4.8 è
\campo{10fg}, il massimo del vento a \SI{10}{m}, con tempo di riferimento $H-1$, intervallo
dichiarato 0--1 e \campo{stepType=max} (\cref{cap:raffica}).

Il file orario delle 00 UTC del giorno D viene per intero, campi accumulati compresi, dalla
trentesima ora della run partita alle 18 UTC di D$-2$; i campi accumulati sono la differenza fra la
trentesima e la ventinovesima ora \parencite{issue71,issue72}. I campi accumulati del file orario
delle 01 UTC di D sono la differenza fra la settima e la sesta ora della run partita alle 18 UTC di
D$-1$. I 24 istanti di validità dalle 01 UTC di D alle 00 UTC di D$+1$ vengono tutti da questa run, e
la somma dei loro campi accumulati è il totale della run dalla sesta alla trentesima ora. Quando il
file wrfout della trentesima ora di una run manca alla sorgente, ogni messaggio del file orario
delle 00 UTC successivo è fatto di punti mancanti; al 14 ottobre 2026 non succede in nessun file
orario dell'archivio.

Poiché i campi istantanei hanno tempo di riferimento $H$ e i campi con template 4.8 hanno $H-1$, un
file orario contiene due tempi di riferimento, e l'indice dei file orari va costruito sull'istante
di validità (\campo{validityDate}, \campo{validityTime}), mai sul tempo di riferimento.
Per distinguere un campo accumulato da uno istantaneo servono \campo{stepType} (o
\campo{typeOfStatisticalProcessing}) e \campo{stepRange}. La chiave
\campo{generatingProcessIdentifier} vale 128, il valore predefinito di eccodes per GRIB2, in tutti
i 246 messaggi, e non separa i campi accumulati da quelli istantanei \parencite{issue53}.

\section{Il cambio di run}
\label{sec:cambio-run}

Il \emph{cambio di run} è il passaggio, una volta al giorno, fra l'istante di validità delle 00 UTC
di D, trentesima ora della run partita alle 18 UTC di D$-2$, e quello delle 01 UTC di D, settima ora
della run partita alle 18 UTC di D$-1$. Il salto fra i due istanti di validità è la differenza fra
le due run alle 00 UTC di D, trentesima ora della prima e sesta della seconda, più la variazione
dell'ora 00--01 UTC nella seconda. La \cref{tab:cambio-run} lo
confronta con la variazione di un'ora dentro la stessa run, sul dominio intero, su un solo cambio di
run, dalle 23 UTC del 18 gennaio alle 01 UTC del 19 gennaio 2024 \parencite{issue72}.

\begin{table}[htbp]
  \centering
  \small
  \caption{Scarto quadratico medio fra due istanti di validità consecutivi, sul dominio intero, 18 e
  19 gennaio 2024.}
  \label{tab:cambio-run}
  \begin{tabular}{@{}lSSSSS@{}}
    \toprule
    Istanti & {\campo{2t} (\si{K})} & {\campo{sp} (\si{\hPa})} & {modulo del vento a \SI{10}{m} (\si{m/s})} & {\campo{t} a \SI{850}{\hPa} (\si{K})} & {\campo{tcwv} (\si{kg/m^2})} \\
    \midrule
    23--00 UTC, stessa run & 0.66 & 0.50 & 1.13 & 0.48 & 0.89 \\
    00--01 UTC, cambio di run & 1.79 & 1.62 & 1.62 & 0.94 & 1.54 \\
    \bottomrule
  \end{tabular}
\end{table}

Il 19 gennaio 2024 lo scarto quadratico medio di \campo{2t} al cambio di run è \num{2.7} volte
quello di un'ora dentro la stessa run, e la differenza assoluta massima fra i due istanti di
validità è \SI{11.0}{K}, contro \SI{9.1}{K} fra le 23 e le 00 UTC. Una
tendenza o una derivata temporale calcolata attraverso l'intervallo fra le 00 e le 01 UTC contiene
questo salto.
